\chapter{La década venezolana}
\label{ch:venezuela}

\section{El tercer patrón}

En octubre de 2000 Fidel Castro y Hugo Chávez firmaron en Caracas un Convenio
Integral de Cooperación. Venezuela suministraría a Cuba crudo y productos
refinados en condiciones concesionales ---al principio unos 53.000 barriles
diarios, que a lo largo de la década subieron hasta unos 90.000 o 100.000--- y
Cuba pagaría, en parte sustancial, no en dinero sino con los servicios de su
gente: médicos, enfermeros, dentistas, maestros, entrenadores deportivos,
técnicos agrícolas y asesores militares y de inteligencia.\unv{}

Vale la pena detenerse en la estructura de este acuerdo, porque es la tercera vez
que la misma estructura aparece en este libro. En 1934 Estados Unidos garantizó un
precio para el azúcar cubano por encima del mercado mundial y Cuba se organizó en
torno a él. Desde 1972 la Unión Soviética garantizó un precio para el azúcar
cubano y suministró petróleo por debajo del mercado mundial, y Cuba se organizó en
torno a eso. Desde 2000 Venezuela suministró petróleo por debajo del mercado
mundial y compró servicios cubanos por encima de su valor de mercado, y Cuba se
organizó en torno a eso. Cada arreglo fue generoso, cada uno fue político y no
comercial, cada uno rescató a la economía cubana en el momento en que se firmó, y
cada uno eliminó la presión que en otro caso habría forzado la diversificación. El
capítulo~\ref{ch:sovieteconomy} enunció la regla; este capítulo es su tercera
confirmación.

\section{Los médicos como industria exportadora}

El intercambio creció mucho más allá del petróleo. Desde 2003, bajo el programa
Barrio Adentro, Cuba situó personal médico en los barrios y zonas rurales de
Venezuela donde nunca habían ejercido médicos venezolanos ---en el pico, algo así
como entre veinte y treinta mil trabajadores de la salud, con un total de personal
civil cubano en Venezuela que llegó quizá a cuarenta mil---.\unv{} Arreglos
similares, menores y en su mayoría pagados en efectivo, se firmaron con Brasil,
Angola, Argelia, Sudáfrica, Qatar y decenas de otros países, y después de 2013 el
programa Mais Médicos colocó unos 11.000 médicos cubanos en regiones desatendidas
de Brasil.

La economía de esto transformó al país. A comienzos de la década de 2010 las
exportaciones cubanas de servicios profesionales valían del orden de ocho a diez
mil millones de dólares al año, de los cuales los servicios de salud eran la gran
mayoría ---frente a un turismo de unos dos mil quinientos millones y un níquel de
unos mil quinientos--- \citep{onei,mesalago2013}.\unv{} La mayor exportación de
Cuba había dejado de ser una materia prima. Eran sus médicos. Es un resultado
genuinamente notable para un país de once millones, es el dividendo directo de la
inversión descrita en el capítulo~\ref{ch:socialstate}, y es la evidencia más
fuerte que existe a favor de la proposición del capítulo~\ref{ch:talent} de que la
población educada de Cuba es un activo exportable.

La otra mitad del arreglo hay que informarla con la misma llaneza. El gobierno
anfitrión le pagaba a Cuba; Cuba le pagaba al médico. La parte del médico era
típicamente una fracción pequeña del valor del contrato ---en el programa
brasileño, el más transparentemente documentado, Cuba recibía aproximadamente
cuatro veces lo que le trasladaba al facultativo \citep{kirk2015}---.\unv{} Los
médicos en misión estaban sujetos a reglas que no se aceptarían en la mayoría de
los mercados laborales: en varios destinos el jefe de misión les guardaba el
pasaporte; se restringían el movimiento, la asociación y las relaciones con
locales; y el médico que abandonaba la misión quedaba, hasta la reforma migratoria
de 2013, impedido de regresar a Cuba durante ocho años, lo que en la práctica
significaba la separación de sus hijos. Estados Unidos operó de 2006 a 2017 un
programa de parole para profesionales médicos cubanos concebido específicamente
para inducir deserciones, y varios miles lo tomaron.

Ambas descripciones son ciertas a la vez, y el resumen honesto es que Cuba
construyó una industria exportadora real y valiosa y la financió en parte
subpagando y sobrecontrolando a quienes la producían. El capítulo~\ref{ch:life}
sostiene que el programa debe conservarse y que el arreglo laboral debe cambiar, y
que esa no es solo la respuesta ética sino la que maximiza los ingresos, puesto
que la restricción que hoy limita a la industria es la oferta de médicos
dispuestos.

\section{Lo que compró el dinero}

La relación venezolana coincidió con, y pagó, el período políticamente más
confiado que ha tenido el gobierno cubano desde los años ochenta.

En el exterior compró un bloque. El ALBA lo fundaron Cuba y Venezuela en 2004 y
llegó a incluir a Bolivia, Nicaragua, Ecuador y varios Estados caribeños.
Petrocaribe, desde 2005, extendió el petróleo concesional venezolano por el Caribe
con misiones médicas cubanas adjuntas. El aislamiento regional que había sido
política norteamericana desde 1962 sencillamente se disolvió: la Organización de
Estados Americanos revocó en 2009 la suspensión de Cuba de 1962, la Comunidad de
Estados Latinoamericanos y Caribeños se fundó sin Estados Unidos en 2011 y Cuba la
presidió en 2013, y la Posición Común de la Unión Europea de 1996, que condicionaba
las relaciones a la democratización, fue sistemáticamente sorteada por los Estados
miembros y sustituida formalmente en 2016.

En casa compró la Batalla de Ideas. Empezó a finales de 1999 con el caso de Elián
González, un niño de cinco años que sobrevivió al naufragio del bote de su madre y
al que unos parientes retuvieron en Miami contra la voluntad de su padre. Cuba se
movilizó durante siete meses por su regreso, en marchas diarias y con una tribuna
permanente construida frente a la misión norteamericana, y ganó: el niño volvió a
casa en junio de 2000. Fue la última victoria política inequívoca que el gobierno
cubano ganó en sus propios términos, y la maquinaria de movilización que construyó
se volcó después hacia programas domésticos ---un ejército de jóvenes trabajadores
sociales, cursos universitarios impartidos en cada municipio, la televisión como
medio de enseñanza, y una gran cantidad de gasto de capital decidido por campaña de
masas en vez de por plan---.

Una de esas campañas merece párrafo propio, porque la Parte III depende de ella.
Para 2005 la red eléctrica cubana estaba fallando: las termoeléctricas habían
sobrepasado su vida de diseño, las salidas no programadas eran rutinarias, y una
temporada de huracanes podía dejar al país a oscuras. La Revolución Energética de
2006 atacó el problema por los dos extremos a la vez. Por el lado de la demanda, el
Estado sustituyó prácticamente todas las bombillas incandescentes del país ---unos
nueve millones--- por fluorescentes compactas, distribuyó millones de arroceras y
ollas de presión eficientes, y reemplazó los refrigeradores y aires acondicionados
más viejos. Por el lado de la oferta, instaló unos cuatro mil grupos electrógenos
pequeños distribuidos, unos 1.300 megavatios en total, situados por todo el país
en vez de concentrados en unas pocas plantas grandes.\unv{} Las horas de apagón
cayeron abruptamente. El programa fue caro, los generadores quemaban combustible
caro, y la estrategia se dejó decaer más tarde ---pero el diagnóstico era correcto,
y era correcto de una manera que importa ahora: el problema eléctrico cubano se
ataca mejor con capacidad distribuida y reducción de demanda que reconstruyendo
grandes plantas centrales---. El capítulo~\ref{ch:energy} se apoya directamente en
este precedente, con los generadores sustituidos por fotovoltaica y almacenamiento.

\section{Dos choques, y un historial crediticio}

En agosto, septiembre y noviembre de 2008 tres huracanes ---Gustav, Ike y
Paloma--- cruzaron Cuba en diez semanas. Entre los tres dañaron o destruyeron algo
más de medio millón de viviendas y causaron pérdidas que convencionalmente se
estiman en unos diez mil millones de dólares, del orden de una quinta parte del
producto nacional.\unv{} Fue la peor temporada de huracanes de la historia
registrada de Cuba, y llegó en el mismo trimestre que la crisis financiera
mundial.

La respuesta de Cuba a la consiguiente estrechez de liquidez causó un daño
duradero. A partir de finales de 2008 el Estado congeló las cuentas bancarias de
las empresas extranjeras que operaban en Cuba ---dinero que esas empresas habían
ganado y tenían derecho a repatriar, por un monto de algo más de mil millones de
dólares--- y sencillamente dejó de pagar a los proveedores.\unv{} Algunos saldos
se liberaron en los años siguientes; otros no. La medida resolvió un problema de
caja y creó un problema de reputación, y el problema de reputación ha resultado
mucho más duradero. Las empresas extranjeras hoy fijan el precio del riesgo de
contraparte cubano sobre la base de lo que ocurrió en 2009, y ninguna declaración
de intenciones desde La Habana ha cambiado eso. El capítulo~\ref{ch:capital}
sostiene que esto es un historial crediticio, que los historiales crediticios se
reparan con conducta a lo largo de años y no con legislación, y que toda
estrategia de inversión que no empiece por saldar estos atrasos está proponiendo
construir sobre un impago.

\section{El final}

A Hugo Chávez se le diagnosticó cáncer en 2011 y se trató en La Habana. Murió el
5 de marzo de 2013. La economía venezolana entró entonces en el derrumbe que la ha
definido desde entonces, y las condiciones de que Cuba disfrutaba empezaron a
erosionarse, con envíos que caían año tras año a lo largo de la segunda mitad de
la década, y más rápido después de 2016.

Para entonces Cuba había tenido trece años de petróleo subsidiado y un auge
exportador de servicios, y los había usado, como había usado los años soviéticos,
para posponer en vez de reestructurar. La industria azucarera no se reconstruyó.
La producción nacional de alimentos no subió. La dualidad monetaria no se unificó,
aunque la unificación se anunció como política en 2011 y se preparó durante una
década. El sistema energético no se convirtió. Lo que el país tenía en cambio era
una renta externa más, mayor de lo que parecía, y una década en que la presión
para cambiar se había retirado.

\begin{ledger}{la década venezolana, 2000--2013}
\emph{Logrado.} Una industria exportadora construida a partir del sistema
educativo que desplazó al azúcar, el níquel y el turismo juntos, y que puso
médicos cubanos en lugares que ningún otro país cubría ---un servicio real
realmente prestado, y la mejor evidencia disponible para el argumento de la Parte
III---. El fin del aislamiento regional de Cuba, con la suspensión de la OEA
revocada y la CELAC fundada y presidida. La Revolución Energética de 2006, un
diagnóstico correcto y una cura parcialmente correcta, cuya lógica adopta el
capítulo~\ref{ch:energy}. Recuperación de la peor temporada de huracanes
registrada sin víctimas masivas, que es para lo que sirve un sistema de defensa
civil con alcance de cuadra.

\emph{Costo.} El tercer comprador garantizado a precio político en setenta años,
tomado por tercera vez, con el tercer resultado idéntico: ninguna
diversificación, ninguna reestructuración, y un precipicio al final. Una
industria exportadora financiada pagando a sus trabajadores una fracción del
valor de su contrato y, en varios destinos, quedándose con sus pasaportes. La
congelación de las cuentas de empresas extranjeras en 2009, que compró un año de
liquidez al precio de un historial crediticio que todavía se está pagando. Y trece
años en los que cada problema estructural que este libro ha nombrado fue
identificado, anunciado como política y no resuelto.
\end{ledger}
